% Appendix: Labb 2, from labs/lab2/README.md. The guide's headings are sections, Labb 2-1 and
% Labb 2-2 among them, and its tasks \labtask. The figures are the ones the guide shares with
% lectures/L04 (the pin-outs, the input and output circuit, the breadboard and the XOR of four
% NAND gates) and its own wiring drawing.
%
% Deliberate differences: the lecture links are chapter and section references, and the
% simulator's link is written out.
%
% \labprep[title]{F1}{label}: see lab1.tex, which defines the same helper.

\makeatletter
\providecommand{\labprep}[3][]{\prep{#2\if\relax\detokenize{#1}\relax\else\ -- #1\fi}%
  \phantomsection\def\@currentlabel{#2}\label{#3}}
\makeatother

\chapter{Labb 2: Grindnät med IC-kretsar}
\label{app:lab2}

I Labb~1 byggde du logiska funktioner av kontakter. Nu gör du samma sak med \textbf{IC-kretsar}:
små kretsar ur 74-serien, med fyra eller sex grindar i varje, uppkopplade på ett kopplingsdäck.
Labben hör till kapitel~\ref{c5:ch} och har två delar som görs i samma pass:
\begin{itemize}
  \item \textbf{Labb 2-1} (\secref{lab2:1}): grindnät med AND-, OR- och inverterarkretsar, från ett
    minimerat uttryck till en fungerande krets.
  \item \textbf{Labb 2-2} (\secref{lab2:2}): samma sorts nät med enbart NAND-grindar, och en titt
    på NOR.
\end{itemize}

Teorin finns i kapitel~\ref{c4:ch}: Karnaughdiagram i \secref{c4:app:a} och IC-kretsar,
kopplingsdäck och planering i \secref{c4:app:b}. Hur du hittar felet när något inte fungerar står
i \secref{c5:app:a}.

\section{Mål}
\label{lab2:sec:goals}
Efter labben ska du kunna:
\begin{itemize}
  \item ta fram ett minimerat uttryck med ett Karnaughdiagram och realisera det med IC-kretsar;
  \item läsa en benplacering och koppla upp en IC-krets med matning, ingångar och utgång;
  \item planera en uppkoppling med en grindtilldelning och en kopplingslista, och följa planen;
  \item bygga om ett nät till enbart NAND-grindar, och förklara varför det kan spara kretsar;
  \item verifiera ett nät mot dess sanningstabell, och felsöka det systematiskt när det inte
    stämmer.
\end{itemize}

\section{Utrustning}
\label{lab2:sec:equipment}
Varje station har:
\begin{itemize}
  \item ett kopplingsdäck med 5~V-matning, strömbrytare för ingångarna och lysdioder för
    utgångarna;
  \item kretsarna \textbf{74HC08} (AND), \textbf{74HC32} (OR) och \textbf{74HC04} (inverterare) till
    Labb~2-1;
  \item kretsarna \textbf{74HC00} (NAND) och \textbf{74HC02} (NOR) till Labb~2-2;
  \item kopplingstråd, och en multimeter.
\end{itemize}

Titta efter på din station hur strömbrytarna och lysdioderna är anslutna. Om stationen saknar egna
lysdioder kopplar du en lysdiod med ett förkopplingsmotstånd på 330~Ω från utgången till jord. Om
du använder egna tryckknappar i stället för stationens strömbrytare behöver varje ingång ett
pull-down-motstånd. Båda beskrivs i \secref{c4:b4} och~\ref{c4:b5}, och figur~\ref{lab2:fig:io}
visar dem.

\bookfigure{l04_io_circuit}{Tryckknapp med pull-down-motstånd på en grindingång, och lysdiod med
förkopplingsmotstånd på utgången.}{lab2:fig:io}

Benplaceringarna för kretsarna, sedda ovanifrån med hacket åt vänster, finns i
figur~\ref{lab2:fig:hc08}--\ref{lab2:fig:hc02}.

\bookfigure{l04_dip_74hc08}{Benplacering för 74HC08, fyra AND-grindar.}{lab2:fig:hc08}

\bookfigure{l04_dip_74hc32}{Benplacering för 74HC32, fyra OR-grindar.}{lab2:fig:hc32}

\bookfigure{l04_dip_74hc04}{Benplacering för 74HC04, sex inverterare.}{lab2:fig:hc04}

\bookfigure{l04_dip_74hc00}{Benplacering för 74HC00, fyra NAND-grindar.}{lab2:fig:hc00}

\bookfigure{l04_dip_74hc02}{Benplacering för 74HC02, fyra NOR-grindar. Här är ben 1 en
utgång.}{lab2:fig:hc02}

Figur~\ref{lab2:fig:breadboard} visar kopplingsdäcket, med kretsen över mittspåret och matningen
dragen.

\bookfigure{l04_breadboard}{Kopplingsdäcket, med en IC-krets över mittspåret.}{lab2:fig:breadboard}

\section{Säkerhet}
\label{lab2:sec:safety}
\begin{warning}
\begin{itemize}
  \item Koppla med matningen \textbf{frånslagen}, och slå på först när kopplingen är kontrollerad
    mot din kopplingslista.
  \item \textbf{Kontrollera kretsens riktning} innan du slår på: hacket åt vänster, ben 1 nere till
    vänster. En felvänd krets får plus och jord omkastade.
  \item Känn efter med ett finger på kretsen strax efter att du slagit på. \textbf{En krets som blir
    varm är felkopplad}: bryt matningen direkt.
  \item Koppla aldrig 24~V från Labb~1 till kopplingsdäcket. Kretsarna tål högst 7~V.
\end{itemize}
\end{warning}

\section{Förberedelser}
\label{lab2:sec:prep}
Görs \textbf{före} passet, hemma eller i samband med kapitel~\ref{c4:ch}. Utan dem hinner du inte
klart, och de ingår i redovisningen.

\labprep[Sanningstabeller]{F1}{lab2:f1}Skriv den förväntade sanningstabellen för en AND-, en OR-
och en inverterargrind, och för näten i Labb~2-1 \hyperref[lab2:1:u2]{uppgift~2}
och~\hyperref[lab2:1:u3]{3}.

\labprep[Karnaughdiagram]{F2}{lab2:f2}Ta fram ett minimerat uttryck för Labb~2-1
\hyperref[lab2:1:u4]{uppgift~4} med ett Karnaughdiagram, och rita grindnätet.

\labprep[NAND-nät]{F3}{lab2:f3}Rita NOT, AND och OR byggda av enbart NAND-grindar (Labb~2-2
\hyperref[lab2:2:u2]{uppgift~2}), och nätet för \code{X = AB + C} med enbart NAND (Labb~2-2
\hyperref[lab2:2:u3]{uppgift~3}). Kontrollera varje nät med De Morgans lagar eller genom att gå
igenom alla rader.

\labprep[Simulera]{F4}{lab2:f4}Bygg varje nät i CircuitVerse,
\url{https://circuitverse.org/simulator}, och kontrollera det mot sanningstabellen, rad för rad.

\labprep[Planera]{F5}{lab2:f5}Gör en \textbf{grindtilldelning} och en \textbf{kopplingslista} för
varje uppgift där du bygger ett nät: Labb~2-1 uppgift~2, 3 och~4 och Labb~2-2 uppgift~3. Använd
formen i \secref{c4:b7}, med bennummer på varje rad.

\filbreak
\section{Labb 2-1}
\label{lab2:1}
AND-, OR- och inverterarkretsar: 74HC08, 74HC32 och 74HC04.

Arbetsgången är densamma i varje uppgift, och det är värt att hålla fast vid den:
\begin{enumerate}
  \item Koppla matningen till varje krets först, ben 14 till +5~V och ben 7 till jord.
  \item Koppla resten enligt kopplingslistan, och bocka av varje rad.
  \item Koppla ingångarna på oanvända grindar till jord.
  \item Slå på, och gå igenom \textbf{alla} kombinationer av ingångarna. Fyll i tabellen:
    förväntat värde från förberedelserna, och uppmätt värde.
  \item Om en rad inte stämmer: felsök enligt \secref{c5:app:a}, inte genom att flytta sladdar på
    måfå.
\end{enumerate}

\labtask{Uppgift 1}{Verifiera AND, OR och inverterare}\phantomsection\label{lab2:1:u1}
Innan du bygger något med en krets ska du veta att den fungerar och att du läser benplaceringen
rätt.
\begin{parts}
  \item Sätt 74HC08 på kopplingsdäcket och koppla matningen. Koppla grind 1 (ben 1, 2 och 3) till
    två strömbrytare och en lysdiod. Koppla ingångarna på grind 2--4 till jord.
  \item Gå igenom alla fyra kombinationerna och fyll i tabellen.
  \item Byt 74HC08 mot 74HC32 \textbf{utan att flytta en enda sladd}, och upprepa. Förklara varför
    det fungerar.
  \item Koppla en inverterare i 74HC04 (ben 1 in, ben 2 ut) och kontrollera båda raderna.
  \item Mät med multimetern spänningen på en utgång som är 1 och en som är 0, med lysdioden
    ansluten. Skriv upp värdena; de används i kontrolluppgiften, övning~\ref{c5:ex:10}.
\end{parts}

\begin{filltable}{@{}ccYYYY@{}}
  \thead{A} & \thead{B} & \thead{AND, förväntat} & \thead{AND, mätt} & \thead{OR, förväntat} &
    \thead{OR, mätt} \\
  \midrule
  0 & 0 & & & & \fillrule
  0 & 1 & & & & \fillrule
  1 & 0 & & & & \fillrule
  1 & 1 & & & & \\
\end{filltable}

\begin{filltable}{@{}cYY@{}}
  \thead{A} & \thead{NOT, förväntat} & \thead{NOT, mätt} \\
  \midrule
  0 & & \fillrule
  1 & & \\
\end{filltable}

\redovisa{tabellerna och de uppmätta spänningarna.}

\labtask{Uppgift 2}{X = AB + C}\phantomsection\label{lab2:1:u2}
Nätet från \secref{c4:a3}, två grindar i två kretsar. Figur~\ref{lab2:fig:wiring} visar det med
kretsbeteckningar och bennummer, så som kopplingslistan läser det.

\bookfigure{lab2_wiring_and_or}{Grindnätet \code{X = AB + C} med 74HC08 och 74HC32, och
bennumren utskrivna.}{lab2:fig:wiring}

\begin{parts}
  \item Koppla upp nätet enligt din kopplingslista från F5.
  \item Gå igenom alla åtta kombinationerna och fyll i tabellen.
  \item Mät spänningen på U1 ben 3 för raden \code{ABC = 110}. Vilken logisk nivå är det, och
    stämmer det med vad \code{AB} ska vara?
\end{parts}

\begin{filltable}{@{}cccYY@{}}
  \thead{A} & \thead{B} & \thead{C} & \thead{X, förväntat} & \thead{X, mätt} \\
  \midrule
  0 & 0 & 0 & & \fillrule
  0 & 0 & 1 & & \fillrule
  0 & 1 & 0 & & \fillrule
  0 & 1 & 1 & & \fillrule
  1 & 0 & 0 & & \fillrule
  1 & 0 & 1 & & \fillrule
  1 & 1 & 0 & & \fillrule
  1 & 1 & 1 & & \\
\end{filltable}

\redovisa{den fungerande kretsen och tabellen.}

\labtask{Uppgift 3}{Bromsassistenten}\phantomsection\label{lab2:1:u3}
Bromsassistenten från \secref{c2:a10} har fyra ingångar och en utgång, med samma beteckningar som
där:
\begin{itemize}
  \item \code{D}: föraren står på bromspedalen;
  \item \code{S} och \code{R}: avståndssensorn respektive radarn, två detektorer som var och en kan
    se ett hinder;
  \item \code{F}: assistanssystemet rapporterar ett fel;
  \item \code{B}: bilen bromsar.
\end{itemize}

Bilen ska bromsa om föraren bromsar, eller om någon av detektorerna ser ett hinder och systemet
inte rapporterar något fel. Föraren ska alltid kunna bromsa, även när systemet är trasigt.
\begin{parts}
  \item Skriv uttrycket för \code{B} (du tog fram det i kapitel~\ref{c2:ch}; kontrollera det mot
    kravet igen), och rita grindnätet med AND, OR och NOT.
  \item Gör en grindtilldelning. Hur många grindar av varje sort behövs, och vilka kretsar?
  \item Koppla upp nätet och gå igenom alla sexton kombinationerna.
  \item Säkerhetskravet: kontrollera särskilt de åtta raderna där \code{D = 1}. Vad ska \code{B}
    vara på alla dem, oavsett de andra ingångarna?
\end{parts}

\begin{filltable}{@{}ccccYY@{}}
  \thead{D} & \thead{S} & \thead{R} & \thead{F} & \thead{B, förväntat} & \thead{B, mätt} \\
  \midrule
  0 & 0 & 0 & 0 & & \fillrule
  0 & 0 & 0 & 1 & & \fillrule
  0 & 0 & 1 & 0 & & \fillrule
  0 & 0 & 1 & 1 & & \fillrule
  0 & 1 & 0 & 0 & & \fillrule
  0 & 1 & 0 & 1 & & \fillrule
  0 & 1 & 1 & 0 & & \fillrule
  0 & 1 & 1 & 1 & & \fillrule
  1 & 0 & 0 & 0 & & \fillrule
  1 & 0 & 0 & 1 & & \fillrule
  1 & 0 & 1 & 0 & & \fillrule
  1 & 0 & 1 & 1 & & \fillrule
  1 & 1 & 0 & 0 & & \fillrule
  1 & 1 & 0 & 1 & & \fillrule
  1 & 1 & 1 & 0 & & \fillrule
  1 & 1 & 1 & 1 & & \\
\end{filltable}

\redovisa{den fungerande kretsen, tabellen, och svaret på d).}

\labtask{Uppgift 4}{Syntes med Karnaughdiagram}\phantomsection\label{lab2:1:u4}
Ett grindnät med fyra ingångar \code{ABCD} och en utgång \code{X} har sanningstabellen nedan. Du
har tagit fram det minimerade uttrycket i F2.

\begin{narrowtable}{@{}cc@{\hspace{3em}}cc@{}}
  \thead{ABCD} & \thead{X} & \thead{ABCD} & \thead{X} \\
  \midrule
  0000 & 0 & 1000 & 0 \\
  0001 & 1 & 1001 & 1 \\
  0010 & 0 & 1010 & 0 \\
  0011 & 1 & 1011 & 1 \\
  0100 & 0 & 1100 & 1 \\
  0101 & 0 & 1101 & 0 \\
  0110 & 0 & 1110 & 1 \\
  0111 & 0 & 1111 & 0 \\
\end{narrowtable}

\begin{parts}
  \item Visa ditt Karnaughdiagram och ditt minimerade uttryck.
  \item Hur många kretsar behövs? Jämför med vad den direkta avläsningen ur tabellen hade krävt.
  \item Koppla upp nätet och kontrollera alla sexton raderna. Skriv ned förväntat och uppmätt värde
    för varje rad, i en tabell som ovan.
\end{parts}

\redovisa{diagrammet, uttrycket, den fungerande kretsen och tabellen.}

\filbreak
\section{Labb 2-2}
\label{lab2:2}
NAND och NOR: 74HC00 och 74HC02.

I kapitel~\ref{c2:ch} såg du att NAND ensam räcker för att bygga vilken funktion som helst. Här
bevisar du det med riktiga kretsar, och ser vad det kostar och vad det sparar.

\labtask{Uppgift 1}{Verifiera NAND och NOR}\phantomsection\label{lab2:2:u1}
\begin{parts}
  \item Koppla grind 1 i 74HC00 (ingångar ben 1 och 2, utgång ben 3) och kontrollera alla fyra
    raderna.
  \item Byt till 74HC02. \textbf{Titta på benplaceringen först}: i 74HC02 är ben 1 en utgång, och
    ingångarna på grind 1 är ben 2 och 3. Flytta sladdarna därefter, och kontrollera alla fyra
    raderna.
  \item Vad hade hänt om du bara bytt kretsen, som i Labb~2-1 \hyperref[lab2:1:u1]{uppgift~1}~c)?
    Två utgångar hade kopplats mot varandra, eller en strömbrytare mot en utgång. Förklara varför
    det inte är bra.
\end{parts}

\begin{filltable}{@{}ccYYYY@{}}
  \thead{A} & \thead{B} & \thead{NAND, förväntat} & \thead{NAND, mätt} & \thead{NOR, förväntat} &
    \thead{NOR, mätt} \\
  \midrule
  0 & 0 & & & & \fillrule
  0 & 1 & & & & \fillrule
  1 & 0 & & & & \fillrule
  1 & 1 & & & & \\
\end{filltable}

\redovisa{tabellen.}

\labtask{Uppgift 2}{NOT, AND och OR av NAND}\phantomsection\label{lab2:2:u2}
Bygg de tre grundgrindarna av enbart NAND-grindar ur en 74HC00, en i taget, med näten från F3:
\begin{itemize}
  \item \textbf{NOT}: en NAND med båda ingångarna ihopkopplade.
  \item \textbf{AND}: en NAND följd av en NOT.
  \item \textbf{OR}: en NOT på varje ingång, följda av en NAND. Kontrollera med De Morgan:
    \code{(A'B')' = A + B}.
\end{itemize}
\begin{parts}
  \item Hur många NAND-grindar kräver var och en? Ryms alla tre i en 74HC00 samtidigt?
  \item Koppla upp var och en och kontrollera den mot sin sanningstabell.
\end{parts}

\redovisa{varje grind när den fungerar.}

\labtask{Uppgift 3}{X = AB + C med enbart NAND}\phantomsection\label{lab2:2:u3}
Samma funktion som i Labb~2-1 \hyperref[lab2:1:u2]{uppgift~2}, nu enligt \secref{c4:a9}:

\begin{textdrawing}
X = ((AB)' · C')'
\end{textdrawing}

\begin{parts}
  \item Koppla upp nätet i en enda 74HC00.
  \item Gå igenom alla åtta raderna. Tabellen ska vara identisk med den i Labb~2-1
    \hyperref[lab2:1:u2]{uppgift~2}.
  \item Jämför de två lösningarna: antal grindar, antal kretsar och antal sladdar. Vilken hade du
    valt om du skulle bygga hundra exemplar?
\end{parts}

\redovisa{den fungerande kretsen, tabellen och jämförelsen.}

\labtask[Extra]{Uppgift 4}{Enbart NOR, och XOR av fyra NAND}\phantomsection\label{lab2:2:u4}
Frivillig, för den som hinner. Den krävs inte för godkänt.
\begin{parts}
  \item Bygg \code{X = AB + C} med enbart NOR-grindar i en 74HC02, med uttrycket
    \code{X = ((A + C)' + (B + C)')'}. Tänk på benplaceringen.
  \item Bygg en XOR av fyra NAND-grindar i en 74HC00, som i figur~\ref{lab2:fig:xor}.

    \bookfigure{l04_net_xor_nand}{XOR byggd av fyra NAND-grindar.}{lab2:fig:xor}

    \noindent Kontrollera alla fyra raderna, och jämför med sanningstabellen för XOR.
  \item Bromsassistenten från Labb~2-1 \hyperref[lab2:1:u3]{uppgift~3} med enbart NAND-grindar.
    Hur många kretsar behövs?
\end{parts}

\redovisa{det du hunnit.}

\filbreak
\section{Redovisning}
\label{lab2:sec:report}
Bocka av varje uppgift när den är redovisad för läraren. Båda i paret ska kunna förklara
lösningen.

\begin{filltable}{@{}p{6.4em}Lp{5.4em}p{6.4em}@{}}
  \thead{Del} & \thead{Uppgift} & \thead{Krävs för G} & \thead{Redovisad} \\
  \midrule
  Förberedelser & F1--F5 & ja & \fillrule
  Labb 2-1 & 1 -- Verifiera AND, OR och inverterare & ja & \fillrule
  Labb 2-1 & 2 -- X = AB + C & ja & \fillrule
  Labb 2-1 & 3 -- Bromsassistenten & ja & \fillrule
  Labb 2-1 & 4 -- Syntes med Karnaughdiagram & ja & \fillrule
  Labb 2-2 & 1 -- Verifiera NAND och NOR & ja & \fillrule
  Labb 2-2 & 2 -- NOT, AND och OR av NAND & ja & \fillrule
  Labb 2-2 & 3 -- X = AB + C med enbart NAND & ja & \fillrule
  Labb 2-2 & 4 -- Extra & nej & \\
\end{filltable}
